% /root/monodromy, 2026-11-19. Complete root, reduction-descent and obstructions-crystals proof/source audits passed.
\subsection{Every odd degree on the supersingular product}
\label{subsec:ss-all-odd-groups}
\status{Proved and independently checked. The parent, reduction/descent and
obstructions/crystals agents independently audited all three complete
proofs and their prerequisites, and each actually reread the original
transvection-group classification. The conclusion covers every odd
positive degree. The even-degree conclusion is established later in
Theorem~\ref{thm:ss-product-all-full-groups}.}

Put $E:y^2+y=x^3$, $X=E\times E$, and
$L=\mathcal O_E(3O)\boxtimes\mathcal O_E(3O)$ in characteristic two.
We first work over an algebraically closed field.
For $d\ge1$ write $P_d=|L^d|$ and let $h_d$ be the determinant of
relative coherent Frobenius on $P_d$.

\begin{lemma}[All odd norm families and their divisorial exceptions]
\label{lem:ss-all-odd-norm-factors}
For every odd $d\ge3$ there is an etale self-isogeny
$\phi_d:X\to X$ of degree $d^2$, with kernel having injective
coordinate projections, and a norm morphism
\[
 \nu_d:P_1=\mathbf P^8\longrightarrow P_d,
 \qquad \nu_d^*\mathcal O_{P_d}(1)=\mathcal O_{\mathbf P^8}(d^2).
\]
In the tensor monomial basis write
$D(s)=\det(s_{ij})$ and
$Q(s)=s^{[2]t}(J_3\otimes J_3)s$.
Every irreducible factor of $\nu_d^*h_d$ is one of $D,Q$, a
pair-tangency divisor $R_g$, or a triple-identification divisor
$T_{g,\eta}$ for distinct $0,g,\eta\in\ker\phi_d$.
Every $R_g$ and $T_{g,\eta}$ is an irreducible divisor distinct
from $D,Q$, and its generic image lies in $h_d=0$.
The multiplicity of $D$ in $\nu_d^*h_d$ is exactly two.
\end{lemma}

\begin{proof}
Choose $b\in\End(E)$ with $b^\dagger b=bb^\dagger=[d-1]$,
using the explicit quaternion norm construction in
Lemma~\ref{lem:DE-A3-ss-norm}, and put
\[
 \phi_d=\begin{pmatrix}1&b\\-b^\dagger&1\end{pmatrix}.
\]
Then $\phi_d^\dagger\phi_d=\phi_d\phi_d^\dagger=[d]I$.
Its degree is $d^2$, and its kernel is the graph
\[
 K_d=\{(-bP,P):P\in E[d]\}.
\]
The restriction of $b$ to $E[d]$ is invertible, since
$d-1$ is relatively prime to $d$; hence both kernel projections
are injective. All these finite groups are etale since $d$ is odd.
The exact norm-line calculation of
Theorem~\ref{thm:ss-uniform-GL8-block} gives
$\operatorname{Nm}_{\phi_d}L\simeq L^d$:
the numerical classes agree, and their symmetric difference is
geometric two-torsion in $\Pic^0(X)$, which is zero on this
supersingular product. Norms of sections are homogeneous of degree
$d^2$, and the norm of a nonzero section is nonzero. They therefore
give the asserted morphism and pullback line bundle.

We give the incidence details for arbitrary $K_d$, rather than
using the finite orbit calculation for degree three.
For $g\ne0$ the pair-contact conditions at $P,P+g$, with common
tangent $\tau$, impose four linear conditions on $H^0(L)$.
Away from finitely many $P$ the two full first jets have rank six,
and these four conditions have rank four. At the exceptions their
rank is at least three. The main incidence is thus an irreducible
$\mathbf P^4$-bundle over a three-dimensional open of
$X\times\mathbf P^1$, and its exceptional part has dimension at
most six. For triples, both kernel projections are injective, so
the three value tensors are independent everywhere: three distinct
rank-one Segre tensors can be collinear only if one tensor factor
is constant. Hence the triple incidence is an irreducible
$\mathbf P^5$-bundle over $X$, of dimension seven.

Propositions~\ref{prop:ss-uniform-graph-mixed-contact} and
\ref{prop:ss-uniform-tangent-triple} supply smooth sources in
these main incidences with finite translated intersections.
Consequently their incidence maps have finite fibres there and
their images are irreducible divisors. A section invariant under
a nonzero kernel translation cannot be a generic exception: such
a translation must preserve $L$, hence must be three-torsion.
On each elliptic degree-three space a nonzero three-torsion
translation has three distinct projective eigenlines; its tensor
action has eigenspaces of dimension three. Thus invariant curves
lie in a finite union of $\mathbf P^2$'s. If the translation does
not preserve $L$, no divisor in $|L|$ is invariant under it.

Neither collision divisor equals $Q$. Indeed the nondegenerate
Frobenius cubic $Q$ contains no vector subspace of dimension
greater than four: an identically zero restriction to $U$ gives
$U^{[2]t}(J_3\otimes J_3)U=0$, and so $2\dim U\le9$.
The pair and triple coefficient fibres have vector dimensions
five and six. Neither equals $D$, either. At a generic pair,
normalize the two values in each factor to $e_0,e_1$, choose the
common tangent with nonzero first component, and impose
$S_{00}=S_{11}=0$, $S_{01}=S_{10}=1$.
The tangent equations solve for $S_{20},S_{21}$, while the
coefficient of the still free $S_{22}$ in $\det S$ is one.
For a generic triple the first-factor values form a basis, and
the three rows of $S$ may be chosen independently in the three
distinct planes orthogonal to the second-factor values.
Those planes admit three independent row vectors. Thus both
generic coefficient fibres contain invertible matrices.

The source singular-curve discriminant $\Delta$ is irreducible
and has degree $c_2(J^1L)=3L^2=54$, with a generic single ordinary
node. The singular-point incidence is a projective bundle; the
two-point jet bounds give dimension at most six for two singular
points, and the second-jet calculation supplies the ordinary node
and generic separable degree one. These are the same first-jet
and Chern-class calculations written out in
Proposition~\ref{prop:ss-level3-three-witness-reduction}; they
involve only the fixed source linear system and are independent
of the chosen odd isogeny. In particular $\Delta$ is distinct
from the two cubic source factors. It is also distinct from
every collision divisor, since each has a smooth-source witness.
At the generic point of $\Delta$ the source node is not identified
with another source point: imposing a singularity at $P$ and
containment at $P+g$ imposes four conditions, except at finitely
many $P$ where at least three remain, so this extra incidence has
dimension at most six. All other identifications there are nodes.

The source relative determinant is $D^2Q$ by
Proposition~\ref{prop:ss-product-relative-Hasse}.
Outside $\Delta$ and the collision divisors, the norm curve is
nodal with smooth normalization $C_s$. The normalization sequence
adds only a Frobenius-unit graph space. At the generic point of
$\Delta$ the same comparison holds, with its one source node
mapping isomorphically to a target node. Thus $\Delta$ adds no
factor to the pullback determinant. The noninvariant integral
source maps birationally onto its image under the etale isogeny;
all failures of ordinary double identification require a tangency
or a triple. The remaining invariant locus has codimension at
least two. This proves the asserted list of all possible factors.

We verify that every collision divisor does occur in the target
Hasse divisor. For a reduced curve singularity with $r$ branches,
the Frobenius-stable rank of the normalization quotient is $r-1$.
Its branch-constant subspace is $k^r/k$, with bijective Frobenius;
modulo these constants, every branch expansion has positive order,
so a sufficiently high Frobenius power lies in the conductor.
The global normalization sequence therefore gives
\[
 g(Z)-f(Z)=\sum_i\bigl(g(C_i)-f(C_i)\bigr)
                  +\sum_x(\delta_x-r_x+1).
\]
The generic source on each collision divisor is smooth and outside
$D\cup Q$, so its defect is two. A tangency, or at least three
branches in a smooth surface, contributes
$\delta-r+1\ge1$. The target defect is consequently at least
three. The relative coherent dimension subtracts the two ambient
dimensions, whose stable rank is zero, so its Frobenius is not
bijective. This proves the converse inclusion just required.

Finally, at the generic point of $D$ the source is smooth and all
norm identifications are transverse double points. Order the nodes
and their branches by a finite etale base change of this local
family. Its graph term in coherent cohomology is then a finite
free Frobenius-unit subbundle. The odd isogeny identifies the two
ambient coherent cohomology subbundles. Hence the relative target
bundle is an extension of the relative source bundle by that unit
graph bundle, over the local ring itself. Multiplication of
Frobenius determinants in this exact sequence gives vanishing
order exactly two along $D$, as in the source polynomial $D^2Q$.
\end{proof}

\begin{theorem}[Irreducible reduced Hasse divisors in all odd degrees]
\label{thm:ss-all-odd-Hasse}
For every odd $d\ge3$, the relative Hasse divisor on $P_d$ is
geometrically irreducible and reduced. At its geometric generic
point the relative Frobenius has stable rank $9d^2-2$,
one less than its rank $9d^2-1$ off the divisor.
\end{theorem}

\begin{proof}
Use the norm family of the preceding lemma. For each $g\ne0$,
the uniform mixed-contact witness lies on $Q\setminus D$ and
on the main $R_g$. Its target Hasse divisor is smooth there by
Lemma~\ref{lem:ss-mixed-self-tacnode}; the line bundle $L^d$ is
very ample and its polarization is separable because $d$ is odd.
Since $Q$ and $R_g$ are irreducible and their generic images lie
in the target Hasse divisor, they are assigned to the same unique
global component. All $R_g$ are thus assigned to the component
containing the image of $Q$, without requiring any transitivity
of automorphisms on the arbitrary kernel.

For every triple, Proposition~\ref{prop:ss-uniform-tangent-triple}
gives a smooth sharp source with exactly one tangent triple and
otherwise nodes. Lemma~\ref{lem:ss-tangent-triple-bridge} makes
the target Hasse divisor smooth at that point. That point belongs
to its triple divisor and to the tangency divisor of its tangent
pair. It therefore assigns every $T_{g,\eta}$ to the same global
component as $Q$ and all $R_g$. The uniqueness at these smooth
points justifies the assignments even when a pullback factor has
multiplicity.

Any additional target Hasse component would pull back to a
nonzero divisor supported only on $D$, with multiplicity at most
two. Its pullback degree would be at most six. But a nonzero
target hypersurface has pullback degree at least $d^2\ge9$.
The pullback cannot vanish identically, since generic sharp source
curves have nodal norm images and thus maximal relative rank.
This contradiction proves irreducibility. Smoothness at any one
of the mixed-contact witnesses proves multiplicity one.

For the stable-rank assertion, fix a nonzero $g$.
The locus in $R_g$ with globally smooth source, one simple target
tacnode, and otherwise only nodes is nonempty and open, by the
mixed-contact witness and its finite-flat collision analysis.
Since $R_g$ is irreducible and distinct from $D,Q$, this open
contains a sharp source. Its normalization defect is two, and
the single simple tacnode contributes
$\delta-r+1=2-2+1=1$; all nodes contribute zero.
Thus the target has defect exactly three, giving relative stable
rank $9d^2-2$. The rank of a sufficiently high Frobenius iterate
is lower semicontinuous. On the now irreducible Hasse divisor it
is generically at least this value and at most one less than the
full rank. Equality follows.
\end{proof}

\begin{theorem}[Full images in every odd degree on the supersingular product]
\label{thm:ss-all-odd-full-groups}
Let $k$ be any perfect field of characteristic two, and give
$X=E\times E$ the embedding by
$L=\mathcal O_E(3O)\boxtimes\mathcal O_E(3O)$.
For every odd $d\ge1$, the geometric and arithmetic images of
the relative mod-two Hasse--Witt local system on its maximal-rank
smooth degree-$d$ hypersurface-section locus are
\[
 G_{\mathrm{geom},d}=G_{\mathrm{arith},d}
          =\GL_{9d^2-1}(\mathbf F_2).
\]
In particular this local system is geometrically irreducible in
every odd degree.
\end{theorem}

\begin{proof}
Degree one is Theorem~\ref{thm:ss-product-degree-one-full-group}.
Let $d\ge3$ be odd, first over an algebraically closed field,
and put $W_d=H^0(X,L^d)$, $b=\dim W_d=9d^2$.
The universal divisor sequence, Kunneth vanishing
$H^1(X,L^{-d})=0$, and Serre duality identify the relative
coherent bundle on $P_d=\mathbf P(W_d)$ with the kernel of
\[
 W_d^*\otimes\mathcal O_{P_d}(-1)
                  \longrightarrow\mathcal O_{P_d},
\]
namely $\Omega^1_{P_d}$ by the Euler sequence.
It has rank $b-1$, and it carries regular Frobenius on all of
$P_d$, including its singular-divisor parameters. The ambient
$H^1(X,\mathcal O_X)$ has zero Frobenius, so its unit Lang system
is the relative system in question. It is generically bijective
by the exact odd-degree rank theorem.

The preceding Hasse theorem and determinant allocation
(Lemma~\ref{lem:irreducible-hasse-quotients-trivial}) imply that
every nonzero proper invariant subspace of this Lang
representation would have a nonzero trivial quotient.
But Theorem~\ref{thm:ss-product-exact-covectors} gives no nonzero
invariant covectors in odd degree. Hence the geometric
representation is irreducible.

At the generic point of its single noninvertibility divisor,
stable rank drops by exactly one. Lemma~\ref{lem:char2-defect-one-inertia}
therefore shows that the geometric group is generated by
transvections. Explicitly, the strict henselian Lang finite part
fixes a codimension-one subspace pointwise, and the one-dimensional
quotient over $\mathbf F_2$ is also fixed. Each nontrivial inertia
element has rank-one square-zero difference from the identity.
Purity extends the quotient by all such inertia to a finite etale
cover of projective space, which is trivial.

The group contains an actual natural $\GL_8(\mathbf F_2)$ block
fixing a complementary summand, by
Theorem~\ref{thm:ss-uniform-GL8-block}. Thus it preserves no
nondegenerate alternating bilinear form: the natural block has
neither an invariant bilinear form nor an invariant covector,
so it lies in the radical of any invariant form on the whole
representation.
Apply the irreducible transvection-group classification
\cite[Theorem~3.6, printed p.~22]{KantorSeress2001}, whose exact
primary statement was reopened for this application. Over
$\mathbf F_2$ its alternatives are the full special linear
group, a symplectic group, an orthogonal group in its natural
even-dimensional representation, or the indicated symmetric
groups on their deleted permutation modules. All alternatives
except the special linear group preserve a nondegenerate
alternating form: in the orthogonal case use the polar form,
and in the permutation cases use the dot product on the
sum-zero space, taking its constant-line quotient when necessary.
They have just been excluded. Since
$\SL_{b-1}(\mathbf F_2)=\GL_{b-1}(\mathbf F_2)$, the claimed
geometric group follows.

Restriction from the Frobenius-invertibility open to its dense
smooth-divisor open preserves connected components of every finite
etale frame cover, so this is the group on the requested smooth
family. The embedding also gives the full complete linear system
in every degree. Indeed each plane cubic has its section ring
generated in degree one (restrict polynomials on $\mathbf P^2$;
$H^1(\mathbf P^2,\mathcal O(d-3))=0$). The degree-$d$ product
sections are tensors of products of $d$ degree-one sections in
the two factors, hence products of $d$ tensor degree-one sections.
Ambient degree-$d$ forms therefore surject onto $H^0(L^d)$.
The induced ambient parameter map has affine-space fibres after
removing its kernel, so pullback preserves the geometric image.

Finally all geometric frame-cover connectedness assertions persist
under algebraically closed extension. Over an arbitrary perfect
base field the arithmetic image contains the full geometric
general linear group and is itself a subgroup of that same
general linear group. It is therefore equal to it.
\end{proof}

\begin{remark}[Scope of the uniform conclusion]
The theorem closes the entire odd-degree full-image problem for
this fixed supersingular product and its indicated embedding.
Even degrees have generic rank $9d^2-3$ and exactly one invariant
covector. The following first-image sheaf and boundary degree comparison
turn these restrictions into the exact uniform even image in
Theorem~\ref{thm:ss-product-all-full-groups}.
\end{remark}
